\begin{textsample}{POS Dim 2 – vem\_ed\_17 – Score 9.00 – vem\_ed\_17/t074.txt}  \label{ex:f2_pos_123}
QUEM É QUEM Hope Windle é diretora do SUNY COIL Center (State University of New York / EUA), referência mundial em Intercâmbios Virtuais.
A seguir, tem-se uma síntese da \textbf{entrevista} concedida por videoconferência para VEm em 23 de março de 2023.
Sabemos que o SUNY COIL Center está de casa nova, SUNY Oneonta.
Pode nos falar sobre a importância disso?
Agora estamos conectados com estudantes e professores, ao redor do mundo, nacionalmente e dentro da SUNY.
Temos o que chamamos “colaboratório”, algo como uma comunidade de prática, ou incubadora.
É um caminho para compartilhar ideias em nossa comunidade de 153 instituições ao redor do mundo.
Também queremos atrair os alunos quando fazemos apresentações ou eventos, para que tenhamos pessoas localmente também, porque eu estou sempre no Zoom, virtualmente.
Qual foi um dos projetos mais desafiadores com que você trabalhou?
Essa é uma boa questão.
Houve alguns, eu diria que quando fui \textbf{coordenadora} na SUNY Ulster [ 2017-2022 ], \textbf{fizemos} um projeto com um adorável professor de Belarus, ele queria abordar empreendedorismo e turismo.
Mas, durante a colaboração, muitos dos links compartilhados com os alunos eram censurados pelo governo (e não sabíamos disso).
Não era permitido aos estudantes compartilhar vídeos do YouTube e o professor não podia realmente contar o que estava acontecendo.
Os estudantes comentaram: é uma grande ideia empreender projetos turísticos reais, mas não podemos fazer neste país.
Foi interessante reconhecer diferentes formas de governo e como afetam os negócios.
Como a sua experiência de designer \textbf{instrucional} influencia a forma como você dirige o SUNY COIL Center?
Quando se pensa em liderança no COIL, temos Jon Rubin, que foi professor e veio com um ponto de vista de professor: pensar no conteúdo, como \textbf{juntar} os professores e desenvolver os projetos.
Hope Windle, diretora do SUNY COIL Center Quando Mary Lou Forward \textbf{chegou} no SUNY COIL Center [ 2017 ], abrimos para um número muito maior de membros.
E então, como resultado da pandemia para mim, minha vivência em \textbf{design} \textbf{instrucional} me faz querer que a experiência COIL seja flexível para todos, em termos das diferentes formas de fazer, além de fornecer o kit de ferramentas e os recursos para que as pessoas possam visualizar facilmente.
O que você prevê para os projetos COIL nos próximos cinco anos?
Estamos atingindo um momento \textbf{incrível}.
No começo dos projetos COIL [ 2006 ], a pergunta era: como fazemos, como usamos a tecnologia?
Agora, temos questões ligadas à \textbf{esperança} crítica, para que pessoas, especialmente aquelas com dificuldades econômicas, possam redesenhar sua educação.
COIL é um meio para que repensem o que vão fazer de suas vidas.

% matched lemmas: chegar, coordenadora, design, entrevista, esperança, fizemos, incrível, instrucional, juntar
\end{textsample}
